% =====================================================================
\section{Phương pháp đề xuất}
% =====================================================================
\subsection{Tổng quan}

% ---------------------------------------------------------------------
\begin{frame}{Quy trình xây dựng ViNeoBERT: ba giai đoạn}
\vspace{0.4em}
\centering
% Shrink only if the diagram would otherwise run past the text block.
\fitpic{\begin{tikzpicture}[font=\scriptsize,
  % text width (not just minimum width) so long labels wrap instead of
  % widening the node and colliding with its neighbour
  stage/.style={draw=vnNavy,fill=vnLight,rounded corners=2pt,
                text width=2.2cm,minimum height=11mm,align=center,inner sep=3pt},
  src/.style={draw=vnSteel,fill=white,rounded corners=2pt,
              text width=1.85cm,minimum height=8mm,align=center,inner sep=3pt}]

  \node[src] (neo) {NeoBERT\\[-1pt]\tiny kiến trúc + không gian đích};
  \node[src,below=7mm of neo] (vide) {ViDeBERTa\\[-1pt]\tiny mô hình nguồn véc-tơ};

  \node[stage,right=7mm of $(neo)!0.5!(vide)$] (s1)
       {\textbf{Giai đoạn 1}\\ Khởi tạo SALT\\[-1pt]\tiny thân đóng băng};
  \node[stage,right=4mm of s1] (s2)
       {\textbf{Giai đoạn 2}\\ Căn chỉnh\\[-1pt]\tiny thân đóng băng};
  \node[stage,right=4mm of s2] (s3)
       {\textbf{Giai đoạn 3}\\ CPT theo lịch WSD\\[-1pt]\tiny toàn bộ tham số};
  \node[src,right=5mm of s3,fill=vnNavy,text=white,draw=vnNavy] (out)
       {\textbf{ViNeoBERT}};

  \draw[-{Latex[length=2mm]},vnSteel,thick] (neo.east)  -- (s1.west|-neo)  -- (s1.west);
  \draw[-{Latex[length=2mm]},vnSteel,thick] (vide.east) -- (s1.west|-vide) -- (s1.west);
  \draw[-{Latex[length=2mm]},vnSteel,thick] (s1)--(s2);
  \draw[-{Latex[length=2mm]},vnSteel,thick] (s2)--(s3);
  \draw[-{Latex[length=2mm]},vnSteel,thick] (s3)--(out);

  \node[below=2mm of s1,text=vnCite,font=\tiny] {không huấn luyện};
  \node[below=2mm of s2,text=vnCite,font=\tiny] {$\sim$20 triệu token};
  \node[below=2mm of s3,text=vnCite,font=\tiny] {$\sim$5 tỷ token};
\end{tikzpicture}}

\vspace{0.8em}
\begin{columns}[T]
\begin{column}{0.95\textwidth}
\begin{alertblock}{Ba cải tiến của khóa luận so với SALT gốc}
  \small
  \begin{enumerate}\itemsep0pt
    \item \textbf{Mở rộng tập điểm neo} bằng các cặp dịch Anh--Việt khai thác tự động.
    \item \textbf{Khởi tạo riêng lớp đầu ra}, kèm hệ số điều chỉnh theo tần suất
          token và hệ số tỉ lệ $\gamma$.
    \item \textbf{Bổ sung giai đoạn căn chỉnh} với phần thân đóng băng trước khi
          huấn luyện toàn bộ.
  \end{enumerate}
\end{alertblock}
\end{column}
\end{columns}
\end{frame}

% ---------------------------------------------------------------------
\begin{frame}{Chọn ViDeBERTa làm mô hình nguồn véc-tơ}
\vspace{0.3em}
\begin{columns}[T]
\begin{column}{0.56\textwidth}
  \scriptsize
  \begin{tabular}{@{} L{2.3cm} L{2.2cm} L{2.7cm}@{}}
    \toprule
    \rowcolor{vnLight}
    \color{vnNavy}\bfseries Đặc điểm &
    \color{vnNavy}\bfseries PhoBERT &
    \color{vnNavy}\bfseries ViDeBERTa\\
    \midrule
    Token hóa      & BPE            & SentencePiece Unigram\\
    Từ vựng        & 64.000         & $\approx$128.000\\
    Đầu vào        & phải tách từ sẵn & \textbf{văn bản thô}\\
    Điểm log-xác suất & không có    & \textbf{có}\\
    Dữ liệu        & 20GB           & \textbf{$\approx$138GB}\\
    Điểm neo tiềm năng & hạn chế    & \textbf{phong phú}\\
    \bottomrule
  \end{tabular}
  \vspace{3pt}\\
  {\footnotesize\color{vnCite} So sánh theo \citet{Nguyen2020} và \citet{Tran2023}}

  \vspace{0.7em}
  {\footnotesize\color{vnCite} Lưu ý: SALT gốc dùng chính PhoBERT làm mô hình
   nguồn; khóa luận chủ đích đổi sang ViDeBERTa.}
\end{column}

\begin{column}{0.42\textwidth}
  \textbf{Ba lý do}
  \begin{itemize}\itemsep2pt
    \item Xử lý trực tiếp văn bản thô, không phụ thuộc công cụ tách từ
          bên ngoài.
    \item Từ vựng lớn hơn $\Rightarrow$ nhiều ứng viên điểm neo hơn.
    \item Học trên dữ liệu tiếng Việt lớn hơn nhiều $\Rightarrow$ không gian
          véc-tơ mã hóa ngữ nghĩa chính xác hơn.
  \end{itemize}

  \vspace{0.4em}
  \begin{block}{Giảm kích thước từ vựng}
    \small $128.000 \rightarrow \mathbf{30.522}$ token, giữ token đặc biệt và
    chọn theo điểm log-xác suất cao nhất -- để ViNeoBERT có số tham số ngang
    NeoBERT gốc.
  \end{block}
\end{column}
\end{columns}
\end{frame}

% ---------------------------------------------------------------------
\subsection{Tập điểm neo}

\begin{frame}{Tập điểm neo: cầu nối giữa hai không gian}
\vspace{0.3em}
\begin{columns}[T]
\begin{column}{0.44\textwidth}
  \lead{Điểm neo là gì?}

  Cặp token \emph{tương ứng ngữ nghĩa} giữa từ vựng tiếng Việt và tiếng Anh --
  mỗi cặp cho ta một điểm \emph{đã biết} của ánh xạ giữa hai không gian.

  \vspace{0.4em}
  \textbf{Hai vai trò}
  \begin{itemize}\itemsep1pt
    \item Token neo được khởi tạo bằng cách \emph{sao chép} véc-tơ NeoBERT.
    \item Các cặp neo làm \emph{tập giám sát} để học ánh xạ cho những token
          còn lại.
  \end{itemize}

  \vspace{0.3em}
  $\Rightarrow$ Số lượng và chất lượng điểm neo \hl{quyết định} chất lượng
  khởi tạo.
\end{column}

\begin{column}{0.53\textwidth}
  \small
  \begin{tabular}{@{} L{2.2cm} L{2.05cm} L{2.4cm}@{}}
    \toprule
    \rowcolor{vnLight}
    \color{vnNavy}\bfseries Nguồn &
    \color{vnNavy}\bfseries Ví dụ &
    \color{vnNavy}\bfseries Cách xác minh\\
    \midrule
    Trùng mặt chữ & \textit{internet} & dịch máy trả về chính nó\\[2pt]
    Chữ số        & \textit{2024}     & nhận trực tiếp\\[2pt]
    \rowcolor{vnLight}
    \textbf{Cặp dịch} & \textit{nước} $\leftrightarrow$ \textit{water}
                      & \textbf{đóng góp của khóa luận}\\
    \bottomrule
  \end{tabular}

  \vspace{0.6em}
  \begin{alertblock}{Bẫy: đồng tự khác nghĩa}
    \small Việt và Anh cùng dùng chữ Latin, nên trùng chuỗi \emph{chưa chắc}
    trùng nghĩa: \textit{song} (Việt: song song) $\neq$ \textit{song}
    (Anh: bài hát). Khóa luận lọc bằng dịch máy \cit{Junczys2018} --
    chỉ giữ khi bản dịch trùng chính chuỗi gốc.
  \end{alertblock}
\end{column}
\end{columns}
\end{frame}

% ---------------------------------------------------------------------
\begin{frame}{Mở rộng tập điểm neo bằng cặp dịch Anh--Việt}
\framesubtitle{Đóng góp chính ở bước khởi tạo}
\vspace{0.2em}
\begin{columns}[T]
\begin{column}{0.55\textwidth}
  \small
  Trước khi khai thác: lọc để chỉ giữ \textbf{từ tiếng Việt hợp lệ}
  (loại \textit{f, j, w, z}; đúng một cụm nguyên âm; tuân thủ quy tắc
  chính tả \textit{k/gh/ngh}, \textit{q}+\textit{u}\,\dots).

  \vspace{0.5em}
  \begin{tabular}{@{} L{1.0cm} L{3.3cm} L{2.85cm}@{}}
    \toprule
    \rowcolor{vnLight}
    \color{vnNavy}\bfseries Tầng &
    \color{vnNavy}\bfseries Cách khai thác &
    \color{vnNavy}\bfseries Ví dụ\\
    \midrule
    1 & Dịch ngược phải trùng từ gốc & sách $\to$ \textit{book} $\to$ sách\\[2pt]
    2 & Từ ghép, dịch một chiều      & bệnh viện $\to$ \textit{hospital}\\[2pt]
    3 & Từ đơn còn lại, lọc thủ công & bổ sung độ phủ\\
    \bottomrule
  \end{tabular}
\end{column}

\begin{column}{0.43\textwidth}
  \begin{block}{Vì sao dịch ngược lọc được từ đa nghĩa}
    \small Một từ đa nghĩa thường \emph{không} dịch ngược về chính nó, vì bản
    dịch trung gian ứng với một nghĩa khác. Ràng buộc dịch ngược trùng khớp
    do đó tự động giữ lại các cặp tương ứng một--một đáng tin cậy.
  \end{block}

  \vspace{0.3em}
  \textbf{Đánh đổi có chủ đích}
  \begin{itemize}\itemsep1pt
    \item Tầng 1 -- ưu tiên \emph{độ chính xác}.
    \item Tầng 2 -- tăng \emph{số lượng}; từ ghép ít đa nghĩa hơn từ đơn.
    \item Tầng 3 -- vét nốt \emph{độ phủ}.
  \end{itemize}
\end{column}
\end{columns}
\end{frame}

% ---------------------------------------------------------------------
\subsection{Khởi tạo SALT}

\begin{frame}{Khởi tạo mỗi token bằng một ánh xạ tuyến tính cục bộ}
\vspace{0.3em}
\begin{columns}[T]
\begin{column}{0.52\textwidth}
  \lead{Bước 1 -- chọn điểm neo gần nghĩa nhất}

  Đo tương đồng trong không gian véc-tơ tĩnh fastText \cit{Bojanowski2017},
  rồi lọc bằng \emph{sparsemax} \cit{Martins2016}:
  \vspace{-4pt}
  \[ s_{t,a}=\frac{f_t^{\top}f_a}{\|f_t\|\|f_a\|},\qquad
     w_t=\mathrm{sparsemax}(s_t) \]
  \vspace{-8pt}

  {\footnotesize Khác softmax, sparsemax đặt hẳn trọng số của các điểm neo
  xa nghĩa \emph{về 0} -- chỉ những điểm thực sự gần mới tham gia.}

  \vspace{0.5em}
  \textbf{Bước 2 -- học ánh xạ riêng cho token đó}
  \vspace{-4pt}
  \[ X_t=\bigl(A^{\mathrm{src}}_t\bigr)^{+}A^{\mathrm{tgt}}_t,
     \qquad e_t=d_t X_t \]
  \vspace{-8pt}

  {\footnotesize $X_t$ là ánh xạ khớp bình phương tối thiểu, đưa véc-tơ nguồn
  của các điểm neo về sát véc-tơ đích tương ứng.}
\end{column}

\begin{column}{0.45\textwidth}
  \begin{block}{Vì sao \emph{cục bộ} chứ không phải một ánh xạ chung}
    \small Hai không gian được huấn luyện độc lập nên \emph{không đẳng cấu}
    \cit{Ormazabal2019} và có tính bất đẳng hướng \cit{Ethayarajh2019}.
    Một ánh xạ tuyến tính duy nhất cho cả từ vựng là không đủ; mỗi token
    được chiếu dựa trên chính các điểm neo lân cận nó.
  \end{block}

  \vspace{0.4em}
  \textbf{Các trường hợp còn lại}
  \begin{itemize}\itemsep1pt
    \item Ít hơn $k_0=8$ điểm neo $\Rightarrow$ dùng trung bình có trọng số
          của các véc-tơ đích.
    \item Không có véc-tơ fastText $\Rightarrow$ lấy mẫu ngẫu nhiên khớp
          thống kê NeoBERT.
  \end{itemize}
\end{column}
\end{columns}
\end{frame}

% ---------------------------------------------------------------------
\begin{frame}{Chiếu riêng cho cả hai ma trận -- khác biệt so với SALT gốc}
\vspace{0.3em}
\begin{columns}[T]
\begin{column}{0.50\textwidth}
  Vì lớp đầu ra của NeoBERT \emph{không} dùng chung trọng số với lớp véc-tơ
  đầu vào, hai ma trận có \textbf{hình học riêng} và phải được khởi tạo
  độc lập:

  \vspace{0.4em}
  \begin{itemize}\itemsep2pt
    \item \textbf{Ma trận véc-tơ $E$}: véc-tơ đích lấy từ các véc-tơ
          \emph{đầu vào} NeoBERT của điểm neo.
    \item \textbf{Ma trận đầu ra $W_{\mathrm{dec}}$}: véc-tơ đích lấy từ các
          véc-tơ \emph{đầu ra} NeoBERT của cùng tập điểm neo.
  \end{itemize}

  \vspace{0.5em}
  {\footnotesize Sao chép ma trận véc-tơ sang lớp đầu ra -- thói quen kế thừa
  từ các kiến trúc dùng chung trọng số -- gây hại vì hai lớp đảm nhận hai vai
  trò hình học khác nhau \cit{Gao2019}. Thực nghiệm ở phần sau xác nhận
  điều này.}
\end{column}

\begin{column}{0.47\textwidth}
\centering
\fitpic{\begin{tikzpicture}[font=\scriptsize,
  m/.style={draw=vnNavy,rounded corners=2pt,minimum width=2.4cm,
            minimum height=8mm,align=center}]
  \node[m,fill=vnSteel!12] (anchor) {Tập điểm neo\\[-1pt]\tiny dùng chung};
  \node[m,fill=vnLight,below left=9mm and -6mm of anchor] (e) {SALT $\to E$};
  \node[m,fill=vnLight,below right=9mm and -6mm of anchor] (w)
       {SALT $\to W_{\mathrm{dec}}$};
  \node[below=5mm of e,text=vnCite,font=\tiny,align=center]
       {không gian\\véc-tơ đầu vào};
  \node[below=5mm of w,text=vnCite,font=\tiny,align=center]
       {không gian\\véc-tơ đầu ra};
  \draw[-{Latex[length=2mm]},vnSteel,thick] (anchor)--(e);
  \draw[-{Latex[length=2mm]},vnSteel,thick] (anchor)--(w);
\end{tikzpicture}}

\vspace{0.5em}
\begin{alertblock}{Kết quả sớm}
  \small Khởi tạo riêng hai ma trận đạt \hl{67,52} điểm F1, so với
  \textbf{54,09} khi dùng chung trọng số.
\end{alertblock}
\end{column}
\end{columns}
\end{frame}

% ---------------------------------------------------------------------
\begin{frame}{Ba bước tinh chỉnh sau khi chiếu}
\vspace{0.4em}
\begin{columns}[T]
\begin{column}{0.315\textwidth}
  \begin{block}{1. Hiệu chỉnh độ lớn}
    \small
    \[ e_t \leftarrow e_t\cdot\frac{\mu_{\mathrm{neo}}}{\|e_t\|} \]
    Đưa mỗi véc-tơ về \emph{độ lớn trung bình} của NeoBERT, giữ nguyên hướng.

    \vspace{2pt}
    {\footnotesize Cần thiết vì Pre-RMSNorm nhạy với độ lớn đầu vào.}
  \end{block}
\end{column}

\begin{column}{0.335\textwidth}
  \begin{block}{2. Bias theo tần suất token}
    \small
    \[ b_{\mathrm{dec},v}=\log p_v \]
    Khởi tạo bias đầu ra bằng log-tần-suất token tiếng Việt
    \cit{Meister2023}.

    \vspace{2pt}
    {\footnotesize Mô hình khỏi tốn hàng nghìn bước chỉ để học lại phân phối
    tần suất.}
  \end{block}
\end{column}

\begin{column}{0.315\textwidth}
  \begin{block}{3. Thu nhỏ trọng số đầu ra}
    \small
    \[ y=\gamma\,W_{\mathrm{dec}}h+b_{\mathrm{dec}},\;\; \gamma=0{,}1 \]
    Để tín hiệu tần suất \emph{chi phối} giai đoạn đầu.

    \vspace{2pt}
    {\footnotesize Nhân vô hướng nên vẫn giữ nguyên hướng ngữ nghĩa mà SALT
    dựng được.}
  \end{block}
\end{column}
\end{columns}

\vspace{0.5em}
\begin{block}{Kết quả của ba bước: điểm xuất phát tốt hơn hẳn}
  \small
  Mất mát khởi điểm xấp xỉ \textbf{entropy unigram tiếng Việt}
  ($\approx$7,3 nat) thay vì mức đoán ngẫu nhiên $\log|V_t|\approx$
  \textbf{10,3 nat} -- tức mô hình bắt đầu từ chỗ \emph{đã biết} phân phối
  từ vựng tiếng Việt.
\end{block}
\end{frame}

% ---------------------------------------------------------------------
\subsection{Huấn luyện}

\begin{frame}{Giai đoạn 2 và 3: căn chỉnh rồi mới huấn luyện toàn bộ}
\vspace{0.3em}
\begin{columns}[T]
\begin{column}{0.47\textwidth}
  \textbf{Giai đoạn 2 -- căn chỉnh với thân đóng băng}
  \begin{itemize}\itemsep1pt
    \item Đóng băng toàn bộ 28 lớp; chỉ huấn luyện hai ma trận mới trên
          $\sim$20 triệu token.
    \item Để hai ma trận \emph{khớp với nhau và với phần thân} trước khi cho
          phần thân thay đổi \cit{deVries2021}.
    \item Chi phí không đáng kể, nhưng giảm $\approx$26\% perplexity.
  \end{itemize}

  \vspace{0.5em}
  \textbf{Giai đoạn 3 -- tiếp tục tiền huấn luyện}
  \begin{itemize}\itemsep1pt
    \item Toàn bộ tham số, $\sim$5 tỷ token CulturaX tiếng Việt
          \cit{Nguyen2024}.
    \item Khởi động lại tốc độ học khi huấn luyện tiếp \cit{Ibrahim2024}.
  \end{itemize}
\end{column}

\begin{column}{0.49\textwidth}
\centering
\fitpic{\begin{tikzpicture}[font=\scriptsize,x=1cm,y=1cm]
  \draw[->,vnCite] (0,0)--(6.15,0) node[below left=1pt and -2pt,font=\tiny]{bước $s$};
  \draw[->,vnCite] (0,0)--(0,2.1) node[above,font=\tiny]{$\eta$};
  \draw[vnNavy,very thick] (0,0)--(0.9,1.6);
  \draw[vnNavy,very thick] (0.9,1.6)--(4.3,1.6);
  \draw[vnNavy,very thick] (4.3,1.6)
        .. controls (5.1,1.45) and (5.4,0.7) .. (5.9,0);
  \draw[vnSteel,thick,dashed] (2.9,1.6)
        .. controls (3.5,1.35) and (3.7,0.6) .. (4.1,0);
  \node[text=vnNavy,font=\tiny] at (0.45,1.85) {khởi động};
  \node[text=vnNavy,font=\tiny] at (2.6,1.82) {ổn định};
  \node[text=vnNavy,font=\tiny] at (5.7,1.15) {làm nguội};
  \node[text=vnSteel,font=\tiny,align=center] at (3.35,0.42)
       {nhánh\\giữa chừng};
\end{tikzpicture}}

\vspace{2pt}
{\footnotesize\color{vnCite} Lịch tốc độ học Warmup--Stable--Decay
 \citep{Hu2024}}

\vspace{0.3em}
\begin{block}{Vì sao chọn WSD thay vì cosine}
  \small Lịch cosine buộc phải \emph{cam kết trước} tổng số bước. WSD cho phép
  rẽ một nhánh làm nguội ngắn tại bất kỳ mốc nào để lấy mô hình hoàn chỉnh,
  trong khi nhánh chính vẫn tiếp tục -- rất hợp với việc đo chất lượng theo
  từng mốc ngân sách.
\end{block}
\end{column}
\end{columns}
\end{frame}
